 \documentclass[a4paper,10pt]{article}
 \usepackage[utf8]{inputenc}  
 \usepackage[T1]{fontenc}
 \usepackage{graphicx}
 \graphicspath{{images/}} 
 
 \usepackage{wrapfig} 
 \usepackage{amsmath} 
 \usepackage{gensymb}
 \usepackage{fancyhdr}
 \usepackage{lastpage}
 \usepackage{textcomp}
 \usepackage{color}
 \usepackage{courier}
 \usepackage{setspace}
 \usepackage{listings}
 \usepackage{ulem}
 \usepackage{geometry}
 \usepackage{array}
 \geometry{a4paper,
 total = {140mm, 237mm},
 left = 35mm,
 top = 30mm}
 \newcommand{\minus}{\scalebox{0.6}[1.0]{$-$}} 
 \newcommand{\signature}[2]{%
  \par\nobreak\bigskip
  \begin{singlespace}%
  \mbox{}\hfill\begin{tabular}{p{8cm} }
      \rule{8cm}{0.5pt}\newline{}%
        \textbf{#1}\\%
       #2 %
  \end{tabular}%
  \end{singlespace}%
  \medskip%
 }
 \renewcommand{\lstlistingname}{Extrait de code}
 \lstset{
 language=Java,
 basicstyle=\scriptsize\ttfamily, 
 upquote=true,
 aboveskip={1.5\baselineskip},
 columns=fullflexible,
 showstringspaces=false,
 extendedchars=true,
 frame=single,
 breaklines=true,
 showtabs=false,
 showspaces=false,
 showstringspaces=false,
 identifierstyle=\ttfamily,
 keywordstyle=\color[rgb]{0.2,0.2,0.2},
 commentstyle=\color[rgb]{0.133,0.545,0.133},
 stringstyle=\color[rgb]{0.627,0.126,0.941},
 tabsize=5,
 }
 \pagestyle{fancy}
 \fancyhead[R]{Projet final Sigmapi}
 \fancyfoot[C]{\thepage}
 \fancyhead[L]{Loïc Gillioz, Loïc Fracheboud}
 \author{Loïc Gillioz, Loïc Fracheboud}
 \title{Informatique - Projet final \\ \Huge Sigmapi (Patapon-like)}
 \date{13 avril 2016}
 \begin{document}
 \maketitle
 \begin{figure}[!h]
 \centering
 \includegraphics[scale=0.15]{images/icones}
 \end{figure}
 \pagebreak
 
 \section{Description sommaire du jeu}
  Sigmapi est basé sur le jeu sorti sur PSP "Patapon". Ce jeu connu un grand succès et c'est désormais une des références des jeux de rythme.
  \paragraph*{Concept}
  Le jeu place le joueur dans le rôle du dieu des Sigmapis. Les Sigmapis sont des indigènes dont la vocation principale est musicale.
Il suffit de quatre touches pour jouer. Chaque touche est assignée à un tambour particulier. La combinaison de plusieurs tambours permet de créer un rythme assigné à un ordre. Les ordres sont ainsi communiqués aux Sigmapis, selon un rythme en quatre ou cinq temps. Quand un rythme est joué, les Sigmapis le répètent tout en exécutant l'action associée à ce rythme. Le joueur doit alors ruser pour utiliser les rythmes adéquats aux situations, pour terminer le niveau avec un minimum de pertes.
  \paragraph*{Graphismes}
Le thème graphique sera revisité pour faire référence aux logos de la HES-SO Valais, et un style visuel propre à notre jeu sera développé. Le but sera de garder l'univers graphique originel (qui a été largement apprécié à la sortie du jeu) et de le dériver à notre sauce.
  \section{Parlons peu, parlons code}
  Pour créer ce jeu, il nous faudra approcher les thèmes suivants :
 \paragraph{}
 \newcolumntype{M}[1]{>{\raggedright}m{#1}}
\begin{tabular}{lM{8cm}}
    Gestion du temps (timers) & Implémenté \tabularnewline
    Gestion des animations avec des SpriteSheets & Implémenté \tabularnewline
    Gestion des séquences de notes & Implémenté \tabularnewline
    Décors récursifs & En cours \tabularnewline
    Gestion des actions par fonction (p.ex. les archers tirent des flèches, lanciers des lances, les épéistes frappent mais ne lancent rien) & En cours \tabularnewline
    Utilisation de la physique pour les objets balistiques et des collisions & En cours \tabularnewline
    Utilisation des collisions pour appliquer les dégâts aux unités & Non implémenté  \tabularnewline
    Création d'animations sans SpriteSheets (transformations par calcul) & Non implémenté \tabularnewline
    Déplacement de la caméra selon la position des Sigmapis & Non implémenté \tabularnewline
    Adaptation des aptitudes des Sigmapis selon leur niveau et race & Non implémenté \tabularnewline
    Mise en place d'un scénario & Dans un lointain futur \tabularnewline
    
 \end{tabular}
 
 
 \vspace{30pt}
 \begin{wrapfigure}[0]{r}{6cm}
 \vspace{-52pt}
 \centering
 \includegraphics[scale=0.5]{signgillioz}
 \end{wrapfigure}
 \signature{Loïc Gillioz}{Sion, le 21 avril 2016} 
 
 \begin{wrapfigure}[0]{r}{6cm}
 \vspace{-52pt}
 \centering
 \includegraphics[scale=1]{signfracheboud}
 \end{wrapfigure}
 \signature{Loïc Fracheboud}{Sion, le 21 avril 2016} 
 \end{document}